% SPDX-License-Identifier: CC-BY-SA-4.0
% Session: Form parallelism to concurrency
% Exercise: Amdahl's law

\begingroup

\begin{exercice}[Implémentation de la pile en Java]

  On rappelle l'algorithme de pile wait-free inspiré de l'algorithme de Afek, Gafni et Morrison. 
  
  \begin{lstlisting}
    public class WaitFreeStack<T> {
      private final AtomicInteger range = new AtomicInteger(0);
      private final AtomicReferenceArray<T> items = ???;
      public void push(T value) {
        int i = range.getAndIncrement();
        items.set(i, value);
      }
      public T pop() {
        for(int i = range.get()-1; i>=0; i--) {
          T value = items.getAndSet(i, null);
          if(value != null) return value;
        }
        return null;
      }
    }
  \end{lstlisting}

  \begin{question}
  \item Dessinez l'exécution séquentielle de l'algorithme quand un processus exécute \lstinline{push('a'); pop();} à partir d'une pile vide.
  \item Dessinez une exécution concurrente de l'algorithme telle que : un processus $p_1$ exécute \lstinline{push('a');} à partir d'une pile vide,
    un autre processus $p_2$ exécute \lstinline{pop();} qui lui retourne \lstinline{null}, et dans laquelle $p_1$ fait le premier pas de calcul.
  \item Modifiez l'algorithme pour remplacer le tableau infini par une liste chaînée.
    Dans un premier temps, on n'essaiera pas de libérer la mémoire après son utilisation.
  \item Proposez des stratégies pour supprimer les n\oe uds de la liste chaînée correspondant aux éléments qui ont déjà été dépilés.
    Bien identifier les problèmes de concurrence qui peuvent se poser. 
  \end{question}

\end{exercice}

\endgroup
\endinput
